\section{Threats to Validity}
\label{sec:threats}

\myparagraph{Construct validity} The labels come only from reruns, so tests whose flakiness does not reproduce on our machine are labeled stable, and the number of runs matters: fewer runs miss rarely failing tests, and the 20-run round identified only 13 of the 30 flaky tests identified with 100 runs. Test order was not randomized across runs, so order-dependent flakiness is likely underrepresented and our results concern mainly non-order-dependent flakiness. \Den{Confirmar que a ordem dos testes foi a mesma em todas as execuções e informar a configuração de paralelismo.} Environmental disruptions can also create spurious flaky tests; we mitigated them by discarding contaminated rounds and repeating them in isolation. The two rounds also differ in more than the number of runs: the .NET SDK used in the 100-run round was not recorded and may differ from the one resolved in the 20-run round, and Testcontainers was run only 5 times in the second round. Test identity is reconstructed from display names and class hierarchies, and a small number of tests could not be resolved to source code (none of them flaky in the 100-run round).

\myparagraph{Internal validity} The 1:1 undersampling balances the test folds and makes precision and F1 optimistic (LOPO quantifies the difference); correlated tests (for example, three NLog tests that failed together) may inflate the positive class; and the two best models were chosen on the same data on which they are evaluated. Attribute tokens may encode project conventions, and removing them reduces the MCC of the best models by about 0.05--0.07. Tokenization choices, such as which keywords are discarded and whether identifiers are split, affect the vocabulary, and the classifiers were not tuned. The qualitative error analysis was performed by the authors on a few examples and is subject to interpretation bias.

\myparagraph{Conclusion validity} The same 30 positives appear in every draw, so the confidence intervals underestimate uncertainty and, with so few positives, the ordering among classifiers is not reliable. The $p$-values of Table~\ref{tab:marcadores} do not survive a correction for the roughly 48,000 candidate tokens, and the rare markers are exploratory.

\myparagraph{External validity} Project selection was not random: it favored open-source libraries with unit tests and domains in which concurrency is common, and it excluded two projects whose suites depend on external servers, which underrepresents network-related flakiness. None of the projects uses MSTest. The 0.24\% rate and the results may therefore not generalize to applications, UI tests, or proprietary code; moreover, flakiness is concentrated in nine repositories, and the results are specific to the .NET toolchain and to the single Windows machine used, without additional load.
